% ============================================================
\section{Admissibility and Branch Selection}
\label{sec:admissibility}
% ============================================================

Many problems produce more than one algebraic completion.
In Coherence Calculus, admissibility is a \emph{minimum completion} principle
on quadratic forms; no external parameter is required.

\paragraph{Generator versus selector.}
Admissibility does not generate continuation or change. The generator lies at
bedrock in nonretroactive exactification
(Axiom~\ref{ax:nonretroactive-exactification}): if an incidence is nonterminal,
it must continue by strict extension. The role of admissibility begins only
after a realized defect or required lower bound has already been produced. It
then selects a minimum-burden completion among those continuations that are
already compelled on the realized surface.

\begin{definition}[Completion by a boundary term]
\label{def:boundary-completion}\lean{CoherenceCalculus.BoundaryCompletion}
Let $Q_0$ be a positive semidefinite quadratic form on a test space $\cX$.
A \emph{boundary completion} is
\[
Q_\kappa(\Phi) := Q_0(\Phi) + \kappa\,|\Phi(0)|^2,\qquad \kappa\ge 0.
\]
\end{definition}

\begin{definition}[Admissibility (minimum completion)]
\label{def:admissibility-formal}\lean{CoherenceCalculus.AdmissibleCompletion}
Let $D$ be a self-adjoint ``required lower bound'' operator.
A completion $Q$ is \emph{admissible} if:
\begin{enumerate}[label=(\roman*)]
\item \textbf{Completion constraint:} $D + Q \succeq 0$ (bounded below), and
\item \textbf{Minimality:} $Q$ is minimal in PSD order among candidates satisfying (i).
\end{enumerate}
This is the same admissibility notion used by the completion-and-selection
algebra of Section~\ref{sec:constraint-compiler}.
\end{definition}

\begin{remark}[Non-uniqueness]
\lean{CoherenceCalculus.admissible_completion_nonunique_example}
The PSD order is partial, so minimal completions need not be unique.
When multiple minimal completions exist, see
Section~\ref{sec:constraint-compiler} for canonical tie-breakers.
\end{remark}

\subsection{Selection archetype: trace and determinant}

The following minimal algebra captures a large class of ``two-branch'' problems.

\begin{definition}[Local transfer matrix]
\label{def:transfer-matrix}\lean{CoherenceCalculus.TransferMatrix}
Fix a scalar $D$ (a trace/invariant).
Define the $2\times 2$ transfer matrix
\[
M_D := \begin{pmatrix} D & -1\\ 1 & 0 \end{pmatrix},
\qquad \det(M_D)=1,\ \Tr(M_D)=D.
\]
\end{definition}

\begin{lemma}[Eigenvalues]
\label{lem:eigenvalues}\lean{CoherenceCalculus.eigenvalue_equation}
The eigenvalues of $M_D$ satisfy
\[
\lambda^2 - D\lambda + 1 = 0,
\quad\text{so}\quad
\lambda_\pm = \frac{D\pm \sqrt{D^2-4}}{2}.
\]
\end{lemma}

\begin{proof}
Compute $\det(\lambda I - M_D)=\lambda^2-D\lambda+1$.
\end{proof}

\begin{definition}[Ratio iteration]
\label{def:ratio-iteration}\lean{CoherenceCalculus.RatioIteration}
Let $(\Psi_n)$ satisfy $(\Psi_{n+1},\Psi_n)^\top = M_D(\Psi_n,\Psi_{n-1})^\top$.
Define $z_n:=\Psi_n/\Psi_{n-1}$. Then
\[
z_{n+1} = D - \frac{1}{z_n}.
\]
\end{definition}

\begin{proposition}[Branch selection as stability]
\label{prop:branch-selection}\lean{CoherenceCalculus.expanding_branch_selected}
If $D>2$, then $\lambda_+>1$ and $\lambda_-<1$.
Under a minimum-completion admissibility rule, the selected branch is $\lambda_+$.
\end{proposition}

\begin{proof}
For $D>2$, $\sqrt{D^2-4}\in(0,D)$ so $\lambda_+ > 1$ and $\lambda_- \in (0,1)$.
Stability/minimum completion excludes the contracting branch when it corresponds to an ill-posed completion.
\end{proof}

\begin{remark}[Two roots plus selection]
\lean{CoherenceCalculus.expanding_branch_selected}
The ``two roots + selection'' structure is the algebraic core of branch selection:
conservation invariants give the polynomial; admissibility selects the branch.
No external parameter is needed.
\end{remark}

% ============================================================
\subsection{Minimal scalar completion (Rayleigh quotient characterization)}
\label{subsec:scalar-completion}
% ============================================================

This subsection provides a closed-form solution for the simplest admissibility problem: completing along a one-parameter ray.

\begin{theorem}[Minimal scalar completion]
\label{thm:scalar-completion}
\lean{CoherenceCalculus.minimal_scalar_completion}
Let $\cH$ be a finite-dimensional real Hilbert space.
Let $D = D^*$ be self-adjoint (the ``required lower bound'') and let $L = L^* \succeq 0$ be a fixed PSD template.
Consider the ray of completions $\{Q_\nu := \nu L : \nu \ge 0\}$.

Assume the \emph{kernel compatibility} condition: for all $x \in \Ker(L)$, $\langle x, D x \rangle \ge 0$.
(Otherwise, $D + \nu L \succeq 0$ is infeasible for all $\nu$.)

Define the \emph{critical coefficient}
\[
\nu_* := \max_{x \ne 0,\, \langle x, Lx \rangle > 0} \frac{-\langle x, D x \rangle}{\langle x, L x \rangle},
\]
with the convention $\nu_* := 0$ if this maximum is nonpositive.

Then:
\begin{enumerate}[label=(\roman*)]
\item \textbf{Feasibility:} $D + \nu L \succeq 0$ if and only if $\nu \ge \nu_*$.
\item \textbf{Minimality:} $\nu_*$ is the unique minimal admissible scalar completion coefficient.
\item \textbf{Certificate:} Any maximizer $x_*$ achieving $\nu_*$ satisfies $\langle x_*, (D + \nu_* L) x_* \rangle = 0$; this is a \emph{tight mode} certifying that $\nu_*$ is sharp.
\end{enumerate}
\end{theorem}

(Proof in Appendix~\ref{sec:appendix-tower}.)

\begin{remark}[Closed-form minimal completion]
\label{rem:compiler-output}
\lean{CoherenceCalculus.scalar_completion_compiler_output}
Theorem~\ref{thm:scalar-completion} is the closed-form solution for the
simplest admissible cone, namely a one-parameter family.
It produces both a value ($\nu_*$) and a certificate ($x_*$).
The template $L$ is fixed by earlier structural constraints (C1)--(C4); this theorem does not choose $L$.
\end{remark}

\begin{remark}[Generalized eigenvalue formulation]
\label{rem:generalized-eigenvalue}
\lean{CoherenceCalculus.generalized_eigenvalue_formulation}
When $L \succ 0$ (strictly positive definite), $\nu_*$ is the largest eigenvalue of the generalized eigenvalue problem $D x = -\lambda L x$, clamped to be nonnegative.
This connects admissibility to standard spectral theory.
On the current finite Lean surface, this is recorded at the least-feasible-scalar level underlying that spectral interpretation.
\end{remark}
